\documentclass[11pt,letterpaper]{article}

%% ── Packages ─────────────────────────────────────────────────────────────────
\usepackage[margin=1in]{geometry}
\usepackage{amsmath,amssymb,bm,bbm}
\usepackage{graphicx}
\usepackage{booktabs}
\usepackage{multirow}
\usepackage{hyperref}
\usepackage{caption}
\usepackage{subcaption}
\usepackage{xcolor}
\usepackage{enumitem}
\usepackage{microtype}
\usepackage{parskip}
\usepackage{array}

\hypersetup{colorlinks=true, linkcolor=blue, citecolor=blue, urlcolor=blue}

%% ── Macros ───────────────────────────────────────────────────────────────────
\newcommand{\R}{\mathbb{R}}
\newcommand{\W}{\mathbf{W}}
\newcommand{\avec}{\mathbf{a}}
\newcommand{\psivec}{\boldsymbol{\psi}}
\newcommand{\mD}{\mathcal{D}}
\newcommand{\norm}[1]{\left\|#1\right\|}

\begin{document}

\noindent \textbf{ECE 66200: Pattern Recognition} \hfill \textbf{Spring 2026} \\
\vspace{1.5cm}

\begin{center}
    {\Large \textbf{Project 1: From PCA to Neural Networks}} \\
    \vspace{0.5cm}
    {\large [Your Name / Author Name Placeholder]} \\
    \vspace{0.5cm}
    {\large \today}
\end{center}
\vspace{1cm}

%% ── Abstract ─────────────────────────────────────────────────────────────────
\begin{abstract}
This report presents a unified study of the \emph{Decomposed Convolutional
Filter} (DCF) framework across three parts.
\textbf{Part 1} shows that learned CNN filters occupy a remarkably
low-dimensional subspace: PCA truncation requires $K \ge 7$ of 9 components
to preserve accuracy without retraining, yet a single learned DCF atom
($K=1$) achieves 98.96\% on MNIST, and Fourier-Bessel (FB) bases match
this performance with no data-driven atom learning.
\textbf{Part 2} applies DCF to MNIST$\to$SVHN domain adaptation.
Counterintuitively, adapting coefficient matrices $a_k$ while keeping
atoms frozen (Exp~C, 93.3\% at $K=8$) far outperforms adapting atoms
while keeping coefficients frozen (Exp~B, 81.2\%), revealing that atoms
are the \emph{domain-invariant} component for digit domains.
\textbf{Part 3} frames continual deblurring as a coefficient-swapping
problem: shared atoms define a universal deconvolution dictionary, and
per-task coefficient snapshots achieve zero catastrophic forgetting
(backward transfer $= 0$\,dB vs.\ $-13.1$\,dB for the naive baseline).
Together, the three parts demonstrate that the low-rank subspace structure
of CNN filters enables compression, adaptation, and continual learning
within a single unified parameterisation.
\end{abstract}

\tableofcontents
\newpage

%%=============================================================================
%% PART 1
%%=============================================================================

\section{Part 1 --- PCA View of CNN Filters}
\label{sec:part1}

\noindent
\textbf{Background: Decomposed Convolutional Filters (DCF).}
Let $\W \in \R^{c' \times c \times l \times l}$ be a conv layer weight tensor.
Reshaping to $\W' \in \R^{c'c \times l^2}$ and applying the compact SVD gives
\begin{equation}
  \W' = \mathbf{U}\mathbf{S}\mathbf{V}^\top, \qquad
  \mathbf{U}\in\R^{c'c\times r},\;
  \mathbf{S}=\mathrm{diag}(s_1,\ldots,s_r),\;
  \mathbf{V}\in\R^{l^2\times r}.
\end{equation}
Truncating to $K$ components gives the DCF representation \cite{qiu2018dcfnet}:
\begin{equation}
  \W[i,j] \;=\; \sum_{k=1}^{K} a_k[i,j]\,\psivec_k,
  \qquad
  \psivec_k \in \R^{l \times l},\quad a_k \in \R^{c' \times c},
  \label{eq:dcf}
\end{equation}
where $\psivec_k$ are \textbf{filter atoms} and $a_k$ are \textbf{coefficients}.
We use \textbf{mode1} (reconstruct-then-convolve):
$\hat{\W} = \sum_k \avec_k \otimes \psivec_k$, then \texttt{F.conv2d}.
Parameter counts are summarised in Table~\ref{tab:p1_params}.

\begin{table}[h]
\centering
\caption{Parameter count per conv layer ($c$ input, $c'$ output,
         kernel $l{\times}l$, $K$ atoms).}
\label{tab:p1_params}
\begin{tabular}{lcc}
\toprule
Variant & Learnable atoms? & Parameters \\
\midrule
Standard Conv2d     & ---        & $c' c l^2$ \\
DCF (learned atoms) & \checkmark & $K l^2 + c' c K$ \\
FB (fixed atoms)    & $\times$   & $c' c K$ \\
\bottomrule
\end{tabular}
\end{table}

\noindent
\textbf{Architecture.}
AlexNet \cite{krizhevsky2012imagenet} is adapted to $1{\times}32{\times}32$
grayscale MNIST: five $3{\times}3$ conv layers (32, 64, 128, 128, 64 channels)
with BatchNorm, three max-pooling layers, and FC layers
$1024\to512\to256\to10$.
DCF and FB variants replace every \texttt{Conv2d} with \texttt{Conv\_DCF}
(mode1).

% ─────────────────────────────────────────────────────────────────────────────
\subsection{Task 1: PCA Filter Truncation}
\label{sec:p1_task1}

\textbf{Setup.}
The baseline AlexNet is trained to convergence (SGD, lr=0.01, cosine
annealing, 30 epochs, batch size 128).
We then apply SVD to each layer's filter matrix and reconstruct with
$K = 1, \ldots, 9$ components \emph{without any retraining}, measuring
test accuracy at each $K$.

\textbf{Results.}

\begin{figure}[h]
  \centering
  \begin{subfigure}{0.48\textwidth}
    \includegraphics[width=\linewidth]{fig1_accuracy_vs_K.pdf}
    \caption{Test accuracy vs $K$.}
    \label{fig:p1_acc}
  \end{subfigure}\hfill
  \begin{subfigure}{0.48\textwidth}
    \includegraphics[width=\linewidth]{fig2_params_vs_K.pdf}
    \caption{Conv parameter count vs $K$.}
    \label{fig:p1_params}
  \end{subfigure}
  \caption{Task 1 results.
    \textbf{(a)} The PCA curve (blue, no retraining) shows a sharp accuracy
    cliff: $K \le 5$ gives 9.7--26.4\%, $K=7$ recovers 97.0\%, and only
    $K=9$ (full rank) restores the 99.65\% baseline.
    The DCF curve (orange, retrained) contrasts starkly: $K=1$ already
    achieves 98.96\% and $K \ge 2$ matches the baseline, showing that the
    compression penalty vanishes when the basis is jointly optimised.
    \textbf{(b)} Parameter count grows linearly in $K$; any $K < 9$ reduces
    conv parameters below the full-rank baseline (dashed).}
  \label{fig:p1_task1}
\end{figure}

\begin{table}[h]
\centering
\caption{MNIST test accuracy (\%) vs $K$. Baseline Conv2d = 99.65\%.}
\label{tab:p1_acc}
\begin{tabular}{c ccc}
\toprule
$K$ & PCA (no retrain) & DCF (retrained) & FB (retrained) \\
\midrule
1 & 9.74  & 98.96 & 98.34 \\
2 & 9.74  & 99.44 & 99.43 \\
3 & 9.74  & 99.49 & 99.60 \\
4 & 9.74  & 99.57 & 99.59 \\
5 & 26.39 & 99.56 & 99.58 \\
6 & 85.92 & 99.61 & \textbf{99.64} \\
7 & 97.00 & 99.58 & --- \\
8 & 97.21 & 99.64 & --- \\
9 & 99.65 & 99.62 & --- \\
\bottomrule
\end{tabular}
\end{table}

\begin{figure}[h]
  \centering
  \includegraphics[width=0.55\textwidth]{fig7_explained_variance.pdf}
  \caption{Cumulative PCA explained variance per conv layer of the trained
    AlexNet.  Conv1 reaches 83\% at $K=6$ and 95.5\% at $K=8$; deeper
    layers saturate more slowly (73\% at $K=6$).  All layers need $K=9$
    for 100\% variance --- consistent with the accuracy cliff: even 5--17\%
    unexplained variance destroys classification without retraining.}
  \label{fig:p1_ev}
\end{figure}

% ─────────────────────────────────────────────────────────────────────────────
\subsection{Task 2: Fourier-Bessel Bases vs Learned DCF Atoms}
\label{sec:p1_task2}

\textbf{Setup.}
We compare two DCF variants retrained from scratch on MNIST:
\begin{itemize}[leftmargin=*]
  \item \textbf{DCF (learned):} $K \in \{1,\ldots,9\}$,
        \texttt{bases\_grad=True}, random atom initialisation.
  \item \textbf{FB (fixed):} $K \in \{1,\ldots,6\}$, atoms set to the
        6 Fourier-Bessel bases (solutions to the 2-D Helmholtz equation
        on the unit disk), \texttt{bases\_grad=False}.
        Only coefficients $a_k$ are trained.
\end{itemize}

\textbf{Results.}

\begin{figure}[h]
  \centering
  \begin{subfigure}{0.54\textwidth}
    \includegraphics[width=\linewidth]{fig3_fb_dcf_accuracy_vs_K.pdf}
    \caption{Accuracy: FB vs DCF vs PCA.}
    \label{fig:p1_fb_dcf}
  \end{subfigure}\hfill
  \begin{subfigure}{0.43\textwidth}
    \includegraphics[width=\linewidth]{fig4_fb_atoms.pdf}
    \caption{The 6 Fourier-Bessel bases ($3{\times}3$).}
    \label{fig:p1_fb_atoms}
  \end{subfigure}
  \caption{Task 2 results.
    \textbf{(a)} FB and DCF closely track each other across all $K$.
    Both reach $>$99.4\% at $K=2$; FB's plateau at $K=6$ (99.64\%) matches
    DCF's best (99.64\% at $K=8$), confirming that the analytical FB subspace
    spans the same task-relevant directions as data-driven atoms.
    \textbf{(b)} FB bases ordered from DC blob ($k=1$) to fine oscillatory
    patterns ($k=6$), mirroring the frequency content of trained CNN filters.}
  \label{fig:p1_task2}
\end{figure}

\begin{figure}[h]
  \centering
  \includegraphics[width=0.78\textwidth]{fig5_dcf_atoms.pdf}
  \caption{Top PCA atoms from Conv1 (columns of $\mathbf{V}$, SVD of the
    $32{\times}9$ filter matrix).
    Atoms $\psivec_1$--$\psivec_3$ resemble low-frequency edge detectors;
    higher components encode finer spatial structure.  The visual similarity
    to the FB bases in Fig.~\ref{fig:p1_fb_atoms} explains why fixed FB atoms
    are competitive with fully learned DCF atoms.}
  \label{fig:p1_pca_atoms}
\end{figure}

% ─────────────────────────────────────────────────────────────────────────────
\subsection{Task 3: Denoising Autoencoder with DCF Encoder}
\label{sec:p1_task3}

\textbf{Setup.}
A convolutional denoising autoencoder with a three-layer encoder
($1\to32\to64\to64$, BatchNorm, MaxPool) and three ConvTranspose2d decoder
layers restoring $1{\times}28{\times}28$.  Trained on MNIST corrupted by
additive Gaussian noise ($\sigma=0.3$, clipped to $[0,1]$), MSE loss, Adam,
30 epochs.  DCF and FB variants replace the encoder \texttt{Conv2d} layers.
Metric: PSNR (dB) between reconstructed and clean images.

\textbf{Results.}

\begin{figure}[h]
  \centering
  \includegraphics[width=0.55\textwidth]{fig6_psnr_vs_K.pdf}
  \caption{Task 3: Denoising PSNR vs $K$.
    Both DCF (orange) and FB (green) rise steeply from $K=1$ to $K=3$,
    then plateau.  DCF reaches 23.75\,dB at $K=9$; FB reaches 23.71\,dB
    at $K=6$ (its maximum) --- nearly identical ceilings.
    At $K=4$, DCF achieves 23.58\,dB, within 0.17\,dB of the ceiling,
    confirming that denoising filters are also low-rank.}
  \label{fig:p1_psnr}
\end{figure}

% ─────────────────────────────────────────────────────────────────────────────
\subsection{Discussion}
\label{sec:p1_disc}

\textbf{Basis choice matters more than compression ratio.}
At $K=6$, PCA (no retrain) gives 85.9\% while DCF and FB both give
$>$99.6\% --- a 13.7-point gap at \emph{identical parameter counts}.
The difference is that DCF/FB optimise atoms jointly with the task,
whereas PCA inherits atoms from a different optimisation (full-rank training).

\textbf{Low effective rank of the filter manifold.}
The explained-variance curves (Fig.~\ref{fig:p1_ev}) show that $K=5$
captures only 62--74\% of filter variance.  The remaining 26--38\% is
sufficient to destroy classification without retraining, yet insignificant
when the network is trained end-to-end with the compressed parameterisation.

\textbf{FB as a principled alternative.}
FB atoms require no training and are indistinguishable from learned DCF atoms
on both classification (Table~\ref{tab:p1_acc}) and denoising
(Fig.~\ref{fig:p1_psnr}), confirming that digit manifolds align well with
low-frequency spatial modes.

\newpage

%%=============================================================================
%% PART 2
%%=============================================================================

\section{Part 2 --- Domain Adaptation via Selective DCF Freezing}
\label{sec:part2}

\noindent
\textbf{Background.}
A classifier trained on MNIST (handwritten digits, white-on-black) fails on
SVHN (house-number photographs, colour, real-world backgrounds).
SVHN is converted to grayscale by channel averaging:
$I_\text{gray} = (I_R+I_G+I_B)/3$, making both datasets
$1{\times}32{\times}32$ inputs while retaining the distribution gap.
The DCF decomposition $\W[i,j] = \sum_k a_k[i,j]\,\psivec_k$ separates:
\begin{itemize}[leftmargin=*]
  \item $\psivec_k$ (\textbf{atoms}): low-level spatial primitives.
  \item $a_k$ (\textbf{coefficients}): the functional recipe combining them.
\end{itemize}

% ─────────────────────────────────────────────────────────────────────────────
\subsection{Task 1: Baseline --- No Adaptation}
\label{sec:p2_task1}

\textbf{Setup.}
The standard Conv2d AlexNet (same architecture as Part 1) is trained on MNIST
(SGD, lr=0.01, momentum 0.9, cosine annealing, 30 epochs) and evaluated
directly on SVHN without any fine-tuning.

\textbf{Result.}
The MNIST-trained model achieves \textbf{18.2\%} on SVHN vs 99.65\% on
MNIST --- an 81-point drop that quantifies the domain gap and motivates
the adaptation strategies in Task~2.

% ─────────────────────────────────────────────────────────────────────────────
\subsection{Task 2: DCF Adaptation Experiments}
\label{sec:p2_task2}

\textbf{Setup.}
Three experiments use DCF-AlexNet with $K \in \{1,2,4,6,8\}$:
\begin{description}[leftmargin=0pt]
  \item[Exp A (upper bound):] Train DCF from scratch on SVHN.
  \item[Exp B (adapt atoms):] Stage~1: train DCF on MNIST
        (\texttt{bases\_grad=True}, 30 epochs, lr=0.01).
        Stage~2: freeze $a_k$, unfreeze $\psivec_k$, fine-tune on SVHN
        (15 epochs, lr=0.001).
  \item[Exp C (adapt coefficients):] Same Stage~1, then
        freeze $\psivec_k$, unfreeze $a_k$, fine-tune on SVHN.
\end{description}

\textbf{Results.}

\begin{table}[h]
\centering
\caption{SVHN test accuracy (\%) summary at $K=8$.}
\label{tab:p2_main}
\begin{tabular}{llc}
\toprule
Condition & Description & SVHN Acc.\ (\%) \\
\midrule
Task 1: Baseline (Conv) & Trained on MNIST, no fine-tuning & 18.2 \\
\midrule
Exp A ($K=8$) & DCF trained from scratch on SVHN & 94.5 \\
Exp B ($K=8$) & MNIST$\to$SVHN: adapt $\psivec_k$ only & 81.2 \\
Exp C ($K=8$) & MNIST$\to$SVHN: adapt $a_k$ only & \textbf{93.3} \\
Pre-adapt ($K=8$) & MNIST DCF model, no fine-tuning & 20.3 \\
\bottomrule
\end{tabular}
\end{table}

\begin{table}[h]
\centering
\caption{SVHN test accuracy (\%) vs $K$ for all DCF conditions.}
\label{tab:p2_vsK}
\begin{tabular}{c cccc}
\toprule
$K$ & Exp A (upper bound) & Exp B (adapt $\psivec_k$) &
      Exp C (adapt $a_k$) & No adapt \\
\midrule
1 & 86.1 & 68.7 & 68.7 & 15.8 \\
2 & 93.1 & 78.8 & 87.7 & 16.0 \\
4 & 94.5 & 81.0 & 92.4 & 25.4 \\
6 & 94.5 & 81.7 & 92.8 & 21.6 \\
8 & 94.5 & 81.2 & \textbf{93.3} & 20.3 \\
\bottomrule
\end{tabular}
\end{table}

\begin{figure}[h]
  \centering
  \includegraphics[width=0.72\textwidth]{fig_p2_accuracy_vs_K.pdf}
  \caption{SVHN accuracy vs $K$.
    Exp A (red) saturates at $\sim$94.5\% for $K \ge 4$.
    Exp C (orange) closely tracks Exp A, reaching 93.3\% at $K=8$ ---
    within 1.2\% of full retraining --- by updating only $a_k$.
    Exp B (blue) plateaus around 81\%, constrained by its 72-parameter
    per-layer atom budget.
    The pre-adapt baseline (grey) sits at 16--25\%, confirming the domain gap.}
  \label{fig:p2_acc_vs_K}
\end{figure}

\begin{figure}[h]
  \centering
  \includegraphics[width=0.62\textwidth]{fig_p2_summary_bar.pdf}
  \caption{Summary bar chart at $K=8$.
    Exp C closes 98.4\% of the 76-point domain gap; Exp B closes 83\%.}
  \label{fig:p2_summary}
\end{figure}

\begin{figure}[h]
  \centering
  \includegraphics[width=\textwidth]{fig_p2_atoms_comparison_K8.pdf}
  \caption{Conv1 filter atoms ($K=8$) before and after Exp~B adaptation.
    \textbf{Top:} MNIST-trained atoms.
    \textbf{Bottom:} Same atoms after fine-tuning on SVHN.
    Low-to-high-frequency ordering is preserved; higher-frequency atoms
    shift more, reflecting SVHN's richer texture.}
  \label{fig:p2_atoms}
\end{figure}

% ─────────────────────────────────────────────────────────────────────────────
\subsection{Discussion}
\label{sec:p2_disc}

\textbf{The atoms are the domain-invariant component.}
Exp C (93.3\%) $\gg$ Exp B (81.2\%) \emph{reverses the original hypothesis}
that atoms are domain-specific.
For digit domains, MNIST-learned spatial primitives are already suited to SVHN;
what must change is how they are combined ($a_k$).

\textbf{Three reasons:}
\begin{enumerate}[leftmargin=*]
  \item \emph{Shared low-level structure.}  Both domains contain digit strokes
        with similar local edge and frequency content.  Atoms encode this
        shared vocabulary; coefficients encode the domain-specific weighting.

  \item \emph{Parameter budget asymmetry.}  Exp B updates 72 atom params
        per layer (360 total); Exp C updates $\sim$279{,}000 coefficient
        params.  Yet Exp A (full retraining, $\sim$310{,}000) achieves only
        1.2\% more than Exp C, confirming frozen atoms are a genuine
        bottleneck for Exp B --- not just a capacity disadvantage.

  \item \emph{Stable initialisation.}  MNIST-trained $a_k$ encode which
        atomic combinations produce class-discriminative responses.
        Fine-tuning from this point on SVHN is a small adjustment.
        Exp B disrupts this by changing the ingredients while keeping the
        recipes fixed.
\end{enumerate}

\textbf{Comparison with the ACDA literature.}
Wang et al.\ \cite{wang2020dictionary} treat atoms as domain-specific in
object-recognition benchmarks where low-level texture changes dramatically.
For digit domains, both source and target share stroke-level primitives,
so the domain shift is captured in the higher-level coefficient structure.

\newpage

%%=============================================================================
%% PART 3
%%=============================================================================

\section{Part 3 --- Continual Deblurring via DCF Coefficient Swapping}
\label{sec:part3}

% ─────────────────────────────────────────────────────────────────────────────
\subsection{Motivation and Connection to SV-PSF Research}
\label{sec:p3_motivation}

Sequential training causes \emph{catastrophic forgetting}
\cite{french1999catastrophic}: gradient updates for Task~$t$ overwrite
weights learned for earlier tasks.
The na\"{i}ve fix --- one model per task --- incurs memory linear in
the number of tasks.

My personal research on \textbf{image deblurring for spatially varying
PSFs (SV-PSFs)} makes this directly relevant.
A deployed camera encounters many blur conditions over its lifetime
(focus distance, aperture, atmospheric turbulence, sensor aging);
each is a new deblurring task.
A practical system must master new conditions without forgetting old ones,
and without storing a separate CNN per blur type.

% ─────────────────────────────────────────────────────────────────────────────
\subsection{Problem Formulation and Approach}
\label{sec:p3_form}

\textbf{Blur tasks.}

\begin{table}[h]
\centering
\caption{Three sequential synthetic deblurring tasks.}
\label{tab:p3_tasks}
\begin{tabular}{clll}
\toprule
Task $t$ & Blur type & PSF & Parameters \\
\midrule
0 & Gaussian  & $h \propto e^{-\|\mathbf{r}\|^2/2\sigma^2}$ & $\sigma{=}1.5$, $7{\times}7$ \\
1 & Motion    & 1-px line at angle $\alpha$ & $\alpha{=}45^\circ$, $5$\,px \\
2 & Defocus   & $h = \mathbbm{1}[\|\mathbf{r}\| \le R]$ & $R{=}2.5$\,px, $7{\times}7$ \\
\bottomrule
\end{tabular}
\end{table}

\textbf{DCF-CL approach.}
A DCF-based denoising autoencoder separates filters into:
\begin{itemize}[leftmargin=*]
  \item \textbf{Shared atoms $\psivec_k$}: learned on Task~0, frozen permanently.
  \item \textbf{Task-specific coefficients $\avec_k^{(t)}$}: a per-task
        snapshot of how the dictionary is combined.
\end{itemize}
\textbf{Protocol:}
(1)~Train all parameters on Task~0, then freeze $\psivec_k$ and store
$\avec_k^{(0)}$.
(2)~For each Task~$t > 0$: re-initialise $\avec_k^{(t)}$, BN, and decoder;
train only these while $\psivec_k$ stays frozen; store $\avec_k^{(t)}$.
(3)~Inference on Task~$j$: inject $\avec_k^{(j)}$ --- an \emph{exact} restore,
since $\psivec_k$ is unchanged.

\textbf{Forgetting metric.}  Backward transfer (BWT):
\begin{equation}
  \mathrm{BWT} = \frac{1}{T{-}1}\sum_{t=1}^{T-1}
  \bigl(\mathrm{PSNR}_{T-1,\,t-1} - \mathrm{PSNR}^*_{t-1}\bigr),
\end{equation}
where $\mathrm{PSNR}^*_j$ is peak PSNR on Task~$j$ right after its own
training.  Negative BWT $=$ forgetting; DCF-CL gives BWT $= 0$ by construction.

% ─────────────────────────────────────────────────────────────────────────────
\subsection{Experimental Setup}
\label{sec:p3_setup}

Both models use the same autoencoder architecture from Part~1 (three-layer
Conv/DCF encoder, three ConvTranspose decoder layers).
DCF-CL uses $K=6$ atoms.
Both trained with Adam, lr=$10^{-3}$, cosine annealing, 20 epochs per task,
batch size 128.
The \emph{baseline} is the same Conv autoencoder trained sequentially
on Tasks 0$\to$1$\to$2 without any CL mechanism.

% ─────────────────────────────────────────────────────────────────────────────
\subsection{Results}
\label{sec:p3_results}

\begin{figure}[h]
  \centering
  \includegraphics[width=\textwidth]{fig_p3_blur_examples.pdf}
  \caption{Three synthetic blur tasks: PSF kernels (top) and example
    blurred MNIST digits (bottom).
    Task~0 (Gaussian): smooth isotropic blur.
    Task~1 (motion): directional smearing at $45^\circ$.
    Task~2 (defocus): disk kernel with boundary ringing.}
  \label{fig:p3_blur}
\end{figure}

\begin{figure}[h]
  \centering
  \includegraphics[width=\textwidth]{fig_p3_forgetting_curves.pdf}
  \caption{Forgetting curves: PSNR on each task across training phases T0$\to$T1$\to$T2.
    \textbf{Baseline (red):} Task~0 drops from 29.16 to 15.57\,dB after T1
    ($-13.6$\,dB), partially recovering to 21.39\,dB after T2 but remaining
    7.8\,dB below peak.  Task~1 is similarly damaged after T2
    (30.57$\to$12.08\,dB).
    \textbf{DCF-CL (blue):} PSNR is locked at peak for every task across
    all future phases.  BWT: Baseline $= -13.1$\,dB; DCF-CL $= 0.0$\,dB.}
  \label{fig:p3_forgetting}
\end{figure}

\begin{figure}[h]
  \centering
  \includegraphics[width=\textwidth]{fig_p3_psnr_heatmaps.pdf}
  \caption{PSNR matrices (\textbf{rows} = training phase; \textbf{cols} =
    evaluation task).
    \textbf{Left (Baseline):} off-diagonal entries below the diagonal drop
    sharply, quantifying forgetting.
    \textbf{Right (DCF-CL):} each column is constant --- the coefficient
    swap restores the exact task model.}
  \label{fig:p3_heatmaps}
\end{figure}

\begin{table}[h]
\centering
\caption{PSNR (dB) after final training phase (Task~2).}
\label{tab:p3_results}
\begin{tabular}{lcccc}
\toprule
Method & Task 0 (Gaussian) & Task 1 (Motion) & Task 2 (Defocus) & BWT (dB) \\
\midrule
Baseline         & 21.39 & 12.08 & 31.02 & $-13.1$ \\
DCF-CL (swap)    & \textbf{28.89} & \textbf{29.38} & \textbf{29.48} & \textbf{0.0} \\
\bottomrule
\end{tabular}
\end{table}

\begin{figure}[h]
  \centering
  \includegraphics[width=0.60\textwidth]{fig_p3_memory_footprint.pdf}
  \caption{Memory footprint for 3 tasks.
    Baseline: $3 \times 111{,}681 = 335{,}043$ params.
    DCF-CL: $162 + 3 \times 93{,}670 = 281{,}172$ params ($1.19\times$
    fewer).  The modest compression reflects that shared atoms are only
    0.17\% of model parameters; the dominant per-task cost is the decoder
    and BN state.}
  \label{fig:p3_memory}
\end{figure}

% ─────────────────────────────────────────────────────────────────────────────
\subsection{Discussion}
\label{sec:p3_disc}

\textbf{Zero forgetting by construction.}
Injecting $\avec_k^{(j)}$ with frozen $\psivec_k$ algebraically recovers
the Task-$j$ model --- not an approximation.  The PSNR heatmap confirms this:
every column of the DCF-CL matrix is constant to four decimal places.

\textbf{Honest assessment of memory savings.}
With $K=6$, the shared atoms are only 162 of 93,315 parameters (0.17\%).
Per-task storage includes DCF coefficients, BN running stats, and the
full decoder ($\sim$93,670 params), nearly equal to storing a full model.
DCF-CL achieves 1.19$\times$ compression vs the Conv baseline but
\emph{no} compression vs storing three DCF models of the same size.
The advantage would be larger if the decoder were also shared, or if
a deeper/wider encoder were used (atoms become a larger fraction).

\textbf{Connection to SV-PSF research.}
\begin{enumerate}[leftmargin=*]
  \item \emph{New conditions without forgetting:} each new PSF calibration
        state adds a lightweight coefficient snapshot; the shared atom
        dictionary is reused.
  \item \emph{Spatially varying extension:} a spatially varying coefficient
        map $\avec_k^{(x,y)}$ would enable per-patch adaptation while
        sharing global atoms --- a spatially adaptive deconvolution network.
  \item \emph{Online adaptation:} only $\avec_k^{(\text{new})}$ is
        fine-tuned (atoms frozen), converging much faster than full retraining.
\end{enumerate}

\textbf{Limitations.}
Atoms are learned only on Task~0.  If later tasks require frequencies
absent from Gaussian-trained atoms, performance may degrade.
Tasks~1 and~2 (29.38, 29.48\,dB) already exceed Task~0 (28.89\,dB),
suggesting sufficient expressiveness --- but this cannot be guaranteed
for all PSF types.  Future work could grow the atom set incrementally
\cite{miao2021continual}.

\newpage

%%=============================================================================
%% CONCLUSION
%%=============================================================================

\section{Conclusion}
\label{sec:conclusion}

The three parts form a coherent story about the \textbf{low-rank subspace
structure of CNN filters}:

\begin{description}[leftmargin=0pt]
  \item[Part 1 reveals the structure.]
        $K \le 6$ atoms suffice for near-perfect classification and
        denoising when the basis is optimised.
        PCA exposes the structure post-hoc; DCF and FB exploit it during
        training, achieving $>$99.4\% accuracy at $K=2$ where PCA gives
        only 9.7\%.

  \item[Part 2 exploits the structure for transfer.]
        For digit-domain adaptation (MNIST$\to$SVHN), atoms are
        \emph{transferable} and coefficients are \emph{domain-specific}.
        Adapting only $a_k$ (Exp~C) achieves 93.3\%, within 1.2\% of
        full retraining.

  \item[Part 3 exploits the structure for continual learning.]
        Freezing atoms and storing per-task coefficient snapshots yields
        exact zero-forgetting.  The mechanism extends naturally to
        SV-PSF deblurring, where per-condition coefficients enable
        on-device multi-condition deblurring without model proliferation.
\end{description}

\noindent
Together, these results validate that the low-rank subspace structure of
CNN filters is not merely a compression curiosity --- it is the key
architectural principle enabling compression, domain transfer, and
continual learning within a single unified parameterisation.

%%=============================================================================
%% REFERENCES
%%=============================================================================

\begin{thebibliography}{9}

\bibitem{qiu2018dcfnet}
Q.~Qiu, X.~Cheng, and G.~Sapiro.
\newblock {DCFNet}: Deep neural network with decomposed convolutional filters.
\newblock In \emph{ICML}, 2018.

\bibitem{wang2020dictionary}
Z.~Wang, X.~Cheng, G.~Sapiro, and Q.~Qiu.
\newblock A dictionary approach to domain-invariant learning in deep networks.
\newblock In \emph{NeurIPS}, 2020.

\bibitem{miao2021continual}
Z.~Miao, Z.~Wang, W.~Chen, and Q.~Qiu.
\newblock Continual learning with filter atom swapping.
\newblock In \emph{ICLR}, 2022.

\bibitem{krizhevsky2012imagenet}
A.~Krizhevsky, I.~Sutskever, and G.~E. Hinton.
\newblock {ImageNet} classification with deep convolutional neural networks.
\newblock In \emph{NeurIPS}, 2012.

\bibitem{lecun1998gradient}
Y.~LeCun, L.~Bottou, Y.~Bengio, and P.~Haffner.
\newblock Gradient-based learning applied to document recognition.
\newblock \emph{Proceedings of the IEEE}, 86(11):2278--2324, 1998.

\bibitem{netzer2011reading}
Y.~Netzer, T.~Wang, A.~Coates, A.~Brass, B.~Wu, and A.~Y. Ng.
\newblock Reading digits in natural images with unsupervised feature learning.
\newblock \emph{NeurIPS Workshop on Deep Learning and Unsupervised Feature
  Learning}, 2011.

\bibitem{french1999catastrophic}
R.~M. French.
\newblock Catastrophic forgetting in connectionist networks.
\newblock \emph{Trends in Cognitive Sciences}, 3(4):128--135, 1999.

\end{thebibliography}

\end{document}
